\documentclass[aspectratio=169,11pt]{beamer}
% Shared Beamer style for practice contest solution walkthroughs.
\usepackage[utf8]{inputenc}
\usepackage[T1]{fontenc}
\usepackage{lmodern}
\usepackage{amsmath,amssymb}
\usepackage{xcolor}
\usepackage{hyperref}
\usepackage{listings}

\definecolor{CPNavy}{HTML}{172033}
\definecolor{CPBlue}{HTML}{2563EB}
\definecolor{CPGreen}{HTML}{16A34A}
\definecolor{CPOrange}{HTML}{F97316}
\definecolor{CPGray}{HTML}{64748B}
\definecolor{CPLight}{HTML}{F8FAFC}
\definecolor{CPCodeBg}{HTML}{F1F5F9}
\definecolor{CPCodeRule}{HTML}{CBD5E1}

\usetheme{default}
\usefonttheme{professionalfonts}
\setbeamertemplate{navigation symbols}{}
\setbeamersize{text margin left=0.65cm,text margin right=0.65cm}
\setbeamertemplate{itemize item}{\color{CPBlue}$\blacktriangleright$}
\setbeamertemplate{itemize subitem}{\color{CPGray}$\triangleright$}

\setbeamercolor{normal text}{fg=CPNavy,bg=white}
\setbeamercolor{frametitle}{fg=white,bg=CPNavy}
\setbeamercolor{title}{fg=white,bg=CPNavy}
\setbeamercolor{subtitle}{fg=CPLight,bg=CPNavy}
\hypersetup{colorlinks=true,urlcolor=CPBlue}

\setbeamertemplate{frametitle}{%
  \nointerlineskip%
  \begin{beamercolorbox}[wd=\paperwidth,ht=0.88cm,dp=0.22cm,leftskip=0.55cm,rightskip=0.35cm]{frametitle}
    \usebeamerfont{frametitle}\insertframetitle\hfill\small\insertframenumber/\inserttotalframenumber
  \end{beamercolorbox}%
}

\setbeamertemplate{footline}{%
  \leavevmode%
  \hbox{%
    \begin{beamercolorbox}[wd=.58\paperwidth,ht=2.4ex,dp=1ex,leftskip=.5cm]{author in head/foot}%
      \color{CPGray}\insertshorttitle
    \end{beamercolorbox}%
    \begin{beamercolorbox}[wd=.42\paperwidth,ht=2.4ex,dp=1ex,rightskip=.5cm plus1fil]{date in head/foot}%
      \hfill\color{CPGray}\insertshortauthor
    \end{beamercolorbox}%
  }%
}

\setbeamertemplate{title page}{%
  \vspace*{1.25cm}
  \begin{beamercolorbox}[wd=\paperwidth,sep=0.65cm,leftskip=0.15cm,rightskip=0.15cm]{title}
    {\color{white}\Huge\bfseries\inserttitle\par}
    \vspace{0.35cm}
    {\color{CPLight}\Large\insertsubtitle\par}
    \vspace{0.6cm}
    {\color{CPLight}\large\insertauthor\par}
  \end{beamercolorbox}
  \vfill
}

\newcommand{\sectiontag}[1]{\textbf{\color{CPBlue}#1}}
\newenvironment{tightitemize}{\begin{itemize}\small\setlength{\itemsep}{0.18cm}}{\end{itemize}}
\lstdefinestyle{cpcode}{
  language=Python,
  basicstyle=\ttfamily\tiny,
  keywordstyle=\color{CPBlue}\bfseries,
  commentstyle=\color{CPGray}\itshape,
  stringstyle=\color{CPGreen},
  backgroundcolor=\color{CPCodeBg},
  frame=single,
  framerule=0.25pt,
  rulecolor=\color{CPCodeRule},
  showstringspaces=false,
  columns=fullflexible,
  keepspaces=true,
  breaklines=true,
  breakatwhitespace=false,
  tabsize=4,
  xleftmargin=0.08cm,
  xrightmargin=0.08cm,
  aboveskip=0.08cm,
  belowskip=0.08cm
}


\title[texassummers]{Texas Summers}
\subtitle{Day 4 solution walkthrough}
\author[Solution Deck]{Competitive Programming Bootcamp}
\date{}

\begin{document}

\begin{frame}[plain]
  \titlepage
\end{frame}

% BEGIN GENERATED STATEMENT INTERPRETATION
\begin{frame}{Interpret the Word Problem}
\small
\sectiontag{Statement excerpts:}
\begin{tightitemize}
  \item \textit{``minimize the amount they sweat''}
  \item \textit{``Print the path as indexes of the shady spots''}
\end{tightitemize}

\vspace{0.18cm}
\sectiontag{Programming interpretation:}
\begin{tightitemize}
  \item A shady spot resets sun exposure, so each walk segment has independent cost.
  \item The cost of walking between two points is proportional to squared Euclidean distance.
  \item Find the shortest path from dorm to classroom in the complete graph of spots.
\end{tightitemize}
\end{frame}
% END GENERATED STATEMENT INTERPRETATION







\begin{frame}{Problem Model}
\small
\sectiontag{Textbook connection:} CP4 4.4.3 Dijkstra's Algorithm

\vspace{0.25cm}
\sectiontag{Core technique:} Dijkstra on a complete geometric graph

\vspace{0.25cm}
\begin{tightitemize}
  \item Shady spots, dorm, and classroom are vertices.
  \item Every pair of points has an edge weighted by squared Euclidean distance.
  \item The output is the sequence of shady spot indices on a shortest path.
\end{tightitemize}
\end{frame}

\begin{frame}{How to Recognize the Approach}
\begin{tightitemize}
  \item All edge weights are nonnegative, so Dijkstra applies.
  \item The graph is complete, so storing every edge is unnecessary.
  \item Parent pointers are required for path reconstruction.
\end{tightitemize}
\end{frame}

\begin{frame}{Algorithm}
\begin{tightitemize}
  \item Store all shady spots, then append dorm and classroom.
  \item Run Dijkstra from the dorm index.
  \item When visiting u, compute weights from u to every other vertex on demand.
  \item Relax improved distances and record parent[v] = u.
  \item After reaching the classroom, follow parents backward and print shady spots.
\end{tightitemize}
\end{frame}

\begin{frame}{Walkthrough of the Key State}
\begin{tightitemize}
  \item Squared distances preserve ordering compared with actual Euclidean distances.
  \item The priority queue may contain stale entries, so skip d > dist[u].
  \item If no shady spot is on the path, print a single dash.
\end{tightitemize}
\end{frame}

\begin{frame}{Why the Algorithm Is Correct}
\begin{tightitemize}
  \item Dijkstra finalizes vertices in nondecreasing shortest distance because weights are nonnegative.
  \item On-demand edge generation considers the same complete graph as an explicit adjacency list.
  \item Parent pointers are updated exactly when a shortest-distance improvement is found.
\end{tightitemize}
\end{frame}

\begin{frame}{Complexity and Implementation Notes}
\small
\sectiontag{Complexity:} O(V\textasciicircum{}2 log V) time and O(V) memory plus the priority queue.

\vspace{0.3cm}
\begin{tightitemize}
  \item Shady spot output indices are the original 0-based input positions.
  \item The classroom and dorm are not printed as part of the shady path.
\end{tightitemize}
\end{frame}



% BEGIN GENERATED CODE WALKTHROUGH
\begin{frame}[fragile]{Code: Read Shady Spots}
\scriptsize
\sectiontag{Source lines:} \texttt{texassummers.py:46--64}

\vspace{0.08cm}
\begin{columns}[T,totalwidth=\textwidth]
\begin{column}{0.58\textwidth}
\begin{lstlisting}[style=cpcode]
import sys
from heapq import heappush, heappop

def main():
    input = sys.stdin.buffer.readline

    INF = float("inf")

    n = int(input())

    points = []

    for _ in range(n):
        x, y = map(int, input().split())
        points.append((x, y))
\end{lstlisting}
\end{column}
\begin{column}{0.38\textwidth}
\sectiontag{What this chunk does}
\begin{tightitemize}
  \item The first n points are shady spots and keep their input indices.
  \item The coordinate list is the vertex set for the graph.
  \item These are the only intermediate locations that may appear in the output path.
\end{tightitemize}
\end{column}
\end{columns}
\end{frame}

\begin{frame}[fragile]{Code: Add Source and Target}
\scriptsize
\sectiontag{Source lines:} \texttt{texassummers.py:66--78}

\vspace{0.08cm}
\begin{columns}[T,totalwidth=\textwidth]
\begin{column}{0.58\textwidth}
\begin{lstlisting}[style=cpcode]
dorm = tuple(map(int, input().split()))
classroom = tuple(map(int, input().split()))

points.append(dorm)
points.append(classroom)

V = n + 2

s = n
target = n + 1
\end{lstlisting}
\end{column}
\begin{column}{0.38\textwidth}
\sectiontag{What this chunk does}
\begin{tightitemize}
  \item Dorm and classroom are appended after the shady spots.
  \item The dorm index is the Dijkstra source; the classroom index is the target.
  \item V is n + 2 because the graph includes both endpoints.
\end{tightitemize}
\end{column}
\end{columns}
\end{frame}

\begin{frame}[fragile]{Code: Initialize Shortest Paths}
\scriptsize
\sectiontag{Source lines:} \texttt{texassummers.py:79--92}

\vspace{0.08cm}
\begin{columns}[T,totalwidth=\textwidth]
\begin{column}{0.58\textwidth}
\begin{lstlisting}[style=cpcode]
dist = [INF for u in range(V)]
dist[s] = 0

parent = [-1 for u in range(V)]

pq = []
heappush(pq, (0, s))
\end{lstlisting}
\end{column}
\begin{column}{0.38\textwidth}
\sectiontag{What this chunk does}
\begin{tightitemize}
  \item dist holds the best known sweat cost from the dorm.
  \item parent records the predecessor used for path reconstruction.
  \item The priority queue starts with the dorm at cost zero.
\end{tightitemize}
\end{column}
\end{columns}
\end{frame}

\begin{frame}[fragile]{Code: Pop the Next Vertex}
\scriptsize
\sectiontag{Source lines:} \texttt{texassummers.py:93--104}

\vspace{0.08cm}
\begin{columns}[T,totalwidth=\textwidth]
\begin{column}{0.58\textwidth}
\begin{lstlisting}[style=cpcode]
while len(pq) > 0:
    d, u = heappop(pq)  # shortest unvisited u

    if d > dist[u]:
        continue

    if u == target:
        break
\end{lstlisting}
\end{column}
\begin{column}{0.38\textwidth}
\sectiontag{What this chunk does}
\begin{tightitemize}
  \item Dijkstra processes the currently cheapest unsettled vertex.
  \item Old heap entries are skipped if they no longer match dist.
  \item Once the classroom is popped, its shortest path is finalized.
\end{tightitemize}
\end{column}
\end{columns}
\end{frame}

\begin{frame}[fragile]{Code: Relax Complete-Graph Edges}
\scriptsize
\sectiontag{Source lines:} \texttt{texassummers.py:106--137}

\vspace{0.08cm}
\begin{columns}[T,totalwidth=\textwidth]
\begin{column}{0.58\textwidth}
\begin{lstlisting}[style=cpcode]
x1, y1 = points[u]

for v in range(V):

    if u == v:
        continue

    x2, y2 = points[v]

    w = (x1 - x2) * (x1 - x2) + (y1 - y2) * (y1 - y2)

    if dist[u] + w >= dist[v]:
        continue

    dist[v] = dist[u] + w

    parent[v] = u

    heappush(pq, (dist[v], v))
\end{lstlisting}
\end{column}
\begin{column}{0.38\textwidth}
\sectiontag{What this chunk does}
\begin{tightitemize}
  \item The graph is complete, so neighbors are generated on the fly.
  \item Squared distance is the edge weight.
  \item A successful relaxation updates both dist and parent, then queues the improved vertex.
\end{tightitemize}
\end{column}
\end{columns}
\end{frame}

\begin{frame}[fragile]{Code: Reconstruct Shady Spots}
\scriptsize
\sectiontag{Source lines:} \texttt{texassummers.py:139--161}

\vspace{0.08cm}
\begin{columns}[T,totalwidth=\textwidth]
\begin{column}{0.58\textwidth}
\begin{lstlisting}[style=cpcode]
path = []

current = parent[target]

while current != s:

    path.append(current)

    current = parent[current]

path.reverse()
\end{lstlisting}
\end{column}
\begin{column}{0.38\textwidth}
\sectiontag{What this chunk does}
\begin{tightitemize}
  \item The parent chain is followed backward from the classroom.
  \item Only vertices before n are shady spots that should be printed.
  \item The path is reversed to recover walking order from dorm to class.
\end{tightitemize}
\end{column}
\end{columns}
\end{frame}

\begin{frame}[fragile]{Code: Print the Path}
\scriptsize
\sectiontag{Source lines:} \texttt{texassummers.py:163--174}

\vspace{0.08cm}
\begin{columns}[T,totalwidth=\textwidth]
\begin{column}{0.58\textwidth}
\begin{lstlisting}[style=cpcode]
    if len(path) == 0:
        print("-")

    else:
        for spot in path:
            print(spot)

main()
\end{lstlisting}
\end{column}
\begin{column}{0.38\textwidth}
\sectiontag{What this chunk does}
\begin{tightitemize}
  \item A direct route with no shady spots prints a single dash.
  \item Otherwise each shady spot index is printed on its own line.
\end{tightitemize}
\end{column}
\end{columns}
\end{frame}
% END GENERATED CODE WALKTHROUGH

\begin{frame}{Source}
\small
\sectiontag{Solution file:} \texttt{practice\_contests/day\_4/texassummers.py}

\vspace{0.3cm}
\sectiontag{Problem:} \url{https://open.kattis.com/problems/texassummers}

\vspace{0.45cm}
Use this deck with the source code open beside it. The slides explain the reasoning; the code shows the exact input handling and output formatting.
\end{frame}

\end{document}
